\section{风险矩阵}
\begin{riskbox}[核心风险]
本轮最大的风险不是产业方向错误，而是盈利兑现速度追不上股价隐含速度。若 2026Q2/Q3 毛利率、现金流或订单节奏低于当前估值要求，光通信板块会出现“业绩增长但股价下跌”的典型估值消化。完整产业链还要额外关注材料/设备良率、运营商 capex、网络设备利润率和光纤光缆价格周期。
\end{riskbox}

\begin{exhibitbox}[表：风险矩阵]
\centering
\small
\begin{tabularx}{\exhibitboxwidth}{>{\bfseries\raggedright\arraybackslash}p{2.2cm} >{\raggedright\arraybackslash\sloppy\hspace{0pt}}X >{\centering\arraybackslash}p{1.2cm} >{\raggedright\arraybackslash}p{3.0cm}}
\toprule
\textbf{风险} & \textbf{表现} & \textbf{等级} & \textbf{监测指标} \\
\midrule
客户集中 & 海外大客户拉货节奏变化直接影响模块厂收入 & 高 & 前五客户、应收账款、存货、订单口径 \\
ASP 下行 & 800G 放量后价格下降快于成本下降 & 高 & 毛利率、单季度收入/利润剪刀差 \\
1.6T 延迟 & 样品验证变慢，规模出货推迟 & 中高 & 公司调研、资本开支、产品认证节奏 \\
CPO 分工变化 & 系统厂/芯片厂拿走部分模块价值 & 中 & 硅光/CPO 合作模式和光源供应方式 \\
材料/设备卡点 & InP/GaAs/薄膜铌酸锂/硅光晶圆、外延、主动耦合和高速测试良率不及预期 & 中高 & 良率、认证周期、资本开支和设备交付 \\
运营商慢周期 & 光纤光缆、OTN、PON、ROADM 和网络设备项目延期 & 中 & 运营商 capex、招标、海缆/FTTH 项目进度 \\
政策与汇率 & 出口管制、海外生产与汇兑扰动 & 中 & 海外收入占比、美元兑人民币、政策公告 \\
拥挤交易 & 板块换手和主题资金集中导致回撤放大 & 高 & 单日成交额、龙虎榜、基金持仓变化 \\
\bottomrule
\end{tabularx}
\end{exhibitbox}

\section{上修与下修触发}
上修触发包括：中际旭创/新易盛 Q2 净利润达到 Q1 的 1.25 倍以上，1.6T 订单进入批量交付，模块毛利率维持 42\% 以上，上游激光芯片/PLC/AWG/调制器通过更多头部客户认证，长飞/亨通/中天现金流和订单同步改善，中兴/烽火网络设备利润率回升。下修触发包括：Q2 净利润低于 Q1 的 0.85 倍，光模块 ASP 超预期下行，经营现金流持续弱于利润，光纤光缆价格再次下行，运营商 capex 延后，或 CPO 架构使可插拔模块价值量被重估。
